\subsection*{Пример 7}

Исходное уравнение на всей действительной прямой:

\[a(x)y^{\prime\prime}(x)+b(x)y^\prime(x)+c(x)y(x)=g(x),\qquad x\in\mathbb R.\]

Все числа ниже сохраняют числовые параметры исходной задачи. Вспомогательные функции полностью определены последовательно; неопределённых коэффициентов или подбираемых параметров нет.

Для компактной записи аргумент x у вспомогательных функций опущен.

\[p_{1}=\frac{1}{1.6732926527362686 \left(x + 0.7294349309869332\right)^{2} + 1}.\]

\[p_{2}=\cos{\left(0.3956791718106786 x - 1.237485158520438 \right)}.\]

\[p_{3}=p_{1} p_{2}.\]

\[q_{1}=-1.3741713202624033.\]

\[q_{2}=\frac{x + 0.4868704629694469}{\sqrt{\left(x + 0.4868704629694469\right)^{2} + 1}}.\]

\[q_{3}=1.262471318706068 q_{1} - 0.7208722265994816 q_{2}.\]

\[r_{1}=\frac{1}{2.7143004385624576 \left(x + 0.844025409063345\right)^{2} + 1}.\]

\[r_{2}=-0.8694060225797431.\]

\[r_{3}=r_{1} r_{2}.\]

\[h_{1}=\cos{\left(0.38896691402710115 x - 0.765129280976121 \right)}.\]

\[d_{1}=- 0.38896691402710115 \sin{\left(0.38896691402710115 x - 0.765129280976121 \right)}.\]

\[e_{1}=- 0.1512952602077663 \cos{\left(0.38896691402710115 x - 0.765129280976121 \right)}.\]

\[h_{2}=- \sin{\left(0.7994408908200522 x + 2.2415464565670673 \right)}.\]

\[d_{2}=- 0.7994408908200522 \cos{\left(0.7994408908200522 x + 2.2415464565670673 \right)}.\]

\[e_{2}=0.6391057379151586 \sin{\left(0.7994408908200522 x + 2.2415464565670673 \right)}.\]

\[h_{3}=h_{1} h_{2}.\]

\[d_{3}=d_{1} h_{2} + d_{2} h_{1}.\]

\[e_{3}=2 d_{1} d_{2} + e_{1} h_{2} + e_{2} h_{1}.\]

\[a(x)=0.6899054725302279 \tanh{\left(p_{3} \right)} + 1.3877465994017353,\quad b(x)=2.928501252330038 \tanh{\left(q_{3} \right)},\quad c(x)=1.9612207794126544 \tanh{\left(r_{3} \right)} - 3.5275716142968054.\]

\[z=\frac{x-(-0.18597579471337672)}{2.721546805554939},\qquad D=1-z^2,\qquad s=1.0323309132762415.\]

\[W(x)=\begin{cases}\exp(1-D^{-1}),&|z|<1,\\0,&|z|\ge1,\end{cases}\qquad y(x)=sW(x)h_{3}.\]

\[\ell(x)=-\frac{2z}{2.721546805554939D^2},\qquad \ell_1(x)=-\frac{2}{(2.721546805554939)^2D^2}-\frac{8z^2}{(2.721546805554939)^2D^3}.\]

\[g(x)=\begin{cases}s e^{1-D^{-1}}\left\{a(x)e_{3}+[2a(x)\ell(x)+b(x)]d_{3}+[a(x)(\ell(x)^2+\ell_1(x))+b(x)\ell(x)+c(x)]h_{3}\right\},&|z|<1,\\0,&|z|\ge1.\end{cases}\]

Здесь $d_{3}$ и $e_{3}$ заданы явными формулами выше и равны первой и второй производным $h_{3}$. Таким образом, в правой части нет невычисленных производных.

Границы носителя: -2.907522600268316 и 2.5355710108415623. Гладкое продолжение решения и правой части за носителем равно нулю; коэффициенты определены на всей прямой.

Резольвентная постановка: $L=a(x)\partial_x^2+b(x)\partial_x+c_R(x)$, где

\[\lambda=3.5275716142968054,\quad c_R(x)=c(x)+3.5275716142968054,\quad g_R(x)=-g(x),\qquad(\lambda-L)y=g_R.\]

\[\underline a=0.6978411268715072,\quad\overline a=2.077652071931963,\quad |b|\le2.928501252330038,\quad |c_R(x)|\le C=1.9612207794126544,\quad\lambda-C\approx1.566350834884151>0.\]

Оценки $\underline a\le a(x)\le\overline a$, $|b(x)|$ и $C$ глобальны и могут быть неточными; они не названы точными экстремумами.

Ищется ограниченное решение на всей прямой. Отрезок [-1,1] - область сравнения сохранённых численных значений, а не граница определения коэффициентов.

Отдельная совместимая задача Коши на $[-1,1]$ имела бы дополнительные условия

\[y(-1)\approx-0.37559836748434083,\qquad y^\prime(-1)\approx-0.4685038562328828.\]

Эти значения следуют из указанного точного решения. Они не заменяют условие ограниченности в резольвентной задаче; произвольные другие начальные данные определяли бы другую задачу Коши.
